Nous avons utilisé des noms pour désigner des valeurs~: les constantes, les paramètres des
procédures et les compteurs des boucles ``for''. Jusqu'à présent, une fois fixée, la valeur désignée par un tel nom
ne changeait plus. Les {\em variables} permettent au contraire de mémoriser des résultats
intermédiaires dans les calculs et de les modifier au cours de l'exécution du programme.

\section{Variables et affectation}
Une variable\index{Variable} est utilisée pour conserver une valeur dans la mémoire
de l'ordinateur. En \python, une variable est créée par une {\em déclaration de
variable}\index{Déclaration de variable}, qui indique le nom de la variable (qui est un identificateur)
et sa valeur initiale, placée après le signe {\idt =}~:

\vspace{\neglskip}\mbox{\blc  \blk\idt  b\ =\ \mcl True}\par
\vspace{\neglskip}\mbox{\blc  \blk\idt  i\ =\ 0}\par
\vspace{\neglskip}\mbox{\blc  \blk\idt  r\ =\ 0.0}\par
\vspace{\neglskip}\mbox{\blc  \blk\idt  s\ =\ {\cha \char34\char34}}\par

\noindent Cela s'écrit exactement comme les déclarations de constantes que nous avons utilisées
jusqu'à présent~: en \python, une constante n'est rien d'autre qu'une variable dont on ne change
jamais la valeur. C'est \python\ qui décide quelle partie de la mémoire et quel codage sera utilisé pour
représenter la valeur de la variable, selon le type de cette valeur~:
{\idt b} est une variable booléenne, {\idt i} une variable entière, {\idt r} une variable
rationnelle et {\idt s} une variable contenant une chaîne de caractères (ici la chaîne vide). En \python, une variable
a donc toujours une valeur, fixée dès sa création. On peut utiliser cette valeur dans une
expression (comme {\idt 2 * i + 1}) en la notant simplement par le nom de la variable.
Il ne faut pas utiliser une variable avant de l'avoir créée~: si le programme commence par
{\idt i = i + 1} sans avoir d'abord donné une valeur initiale à {\idt i}, l'exécution s'arrête
avec le message `{\tt NameError:\ name 'i' is not defined}' (le nom {\idt i} n'est pas défini).
L'instruction
d'{\em affectation\/} \index{Affectation}permet de changer la
valeur d'une variable. Par exemple l'instruction~:

\vspace{\neglskip}\mbox{\blc  \blk\blk\idt  b\ =\ \mcl False}\par

\noindent affecte la valeur ``faux'' à la variable booléenne {\idt b}. De manière générale, une
instruction d'affectation a la forme suivante~:

\index{Affectation}
\begin{center}
\begin{tabular}{@{}|p{12.5cm}|@{}} \hline
{\em Instruction d'affectation~:}\\ \hline \hline
\begin{minipage}[t]{12.5cm}
{\medskip

\vspace{\neglskip}\mbox{\blc  \blk\blk\idt  nomDeVariable\ =\ \rm expression%
\footnote{\ L'expression peut utiliser l'ancienne valeur de la variable.}}\par

\medskip}
\end{minipage}\\ \hline
\end{tabular}
\end{center}

\noindent Le symbole\ {\idt =}\ se lit ``devient égal à''. Il ne faut pas le confondre avec la
comparaison {\idt ==} (``est égal à'') utilisée dans les conditions des {\mcl if} et des {\mcl while}.
L'exécution de cette instruction consiste à
évaluer l'``expression'' pour calculer sa valeur qui devient celle de la variable nommée
{\idt  nomDeVariable}. Le type de la variable est celui de la valeur qu'elle contient~:

\index{Compatibilité de type}
\begin{center}
\begin{minipage}[t]{11.5cm}
\begin{tabular}{@{}|l|l|l|@{}} \hline
\multicolumn{3}{@{}|l|@{}}{\em Types des valeurs~:} \\ \hline \hline
Exemple de valeur & Type & Nom du type\\ \hline
{\mcl True}, {\mcl False} & {\idt bool} & booléen \\
{\idt 0}, {\idt 1}, {\idt -5} & {\idt int} & entier \\
{\idt 0.0}, {\idt 1.5} & {\idt float} & rationnel \\
{\cha \char34\char34}, {\cha \char34{}abc\char34} & {\idt str} & chaîne de caractères \\
\hline
\end{tabular}
\end{minipage}
\end{center}

\noindent  En \python, le type d'une variable n'est pas fixé une fois pour toutes~:
si {\idt i} est une variable entière, l'affectation {\idt i = 1.5} est
acceptée et {\idt i} devient une variable rationnelle. C'est rarement une bonne idée, car le
lecteur du programme s'y perd et l'erreur n'apparaît que plus tard, par exemple dans
{\idt avf(i)} qui exige un entier. Mieux vaut garder à chaque variable le même type et faire les
{\em conversions}\index{Conversion} explicitement~: par exemple
{\idt r = float(i)} ou {\idt i = int(r)}.

\section{Sauvegarde de l'état du robot}

Une première utilisation des noms est la sauvegarde de l'état du \robot. Imaginons par
exemple que nous voulions concevoir une procédure pour faire un dessin qui nécessite divers
changements de la taille de la grille. Pour que l'appel de
cette procédure dans un programme ne modifie pas la taille de la grille utilisée dans le
programme, il faut sauvegarder cette
taille au début de la procédure puis la restituer à la fin. Il suffit de lui donner un
nom par une {\em déclaration locale}\index{Déclaration locale}
{\idt nom = expression} placée dans le corps de la procédure, comme dans l'exemple suivant~:

{\raggedright
 %\vspace{\baselineskip}
 \vbox to 0pt{\includevspacePostScript[13mm]{Canards}{0.95}{0.95}}%
 \input python/Canards.tex
 }

\noindent Les déclarations {\idt lX = valLgX();\ lY = valLgY()}
calculent les valeurs de {\idt valLgX()} et de {\idt valLgY()} au début de la procédure et
leur donnent les noms {\idt lX} et {\idt lY}, qui ne peuvent être utilisés que dans la suite du
corps de la procédure {\idt canard}. On dit que ce sont des variables {\em locales} à la
procédure~: elles n'existent pas en dehors de son corps.

\section{Compter}
Supposons que nous voulions
compter le nombre d'itérations dans une boucle ``while''. On utilise une variable entière~:

\mbox{\blc  \blk\idt  i\ =\ 0}\par

\noindent Avant d'exécuter la boucle ``while'', la variable doit être initialisée à 0
(elle peut avoir été modifiée depuis sa déclaration)~:

\mbox{\blc  \blk\blk\idt  i\ =\ 0}\par

\noindent Après chaque tour de boucle, la valeur de la variable doit être augmentée de 1~:

\mbox{\blc  \blk\blk\idt  i\ =\ i\ +\ 1}\par

\noindent Supposons que la valeur de la variable {\idt i} soit égale à 5 {\em avant} l'affectation
\mbox{\idt  i\ =\ i\ +\ 1}. La valeur de l'expression {\idt  i\ +\ 1}
est $5 + 1 = 6$ de sorte que la valeur de la variable {\idt i} {\em après} l'affectation sera égale
à 6. Comme cette instruction est très fréquente, \python\ permet de l'abréger en
\mbox{\idt  i\ +=\ 1} (``augmenter {\idt i} de 1''). De même \mbox{\idt  i\ -=\ 1} diminue {\idt i}
de 1 et \mbox{\idt  i\ *=\ 2} le double.

Par exemple, ceci peut nous servir pour déterminer la taille de la fenêtre dans laquelle le \robot\
se déplace sur l'écran de l'ordinateur~:
\input python/CalculerTailleFenetreRobot.tex

\noindent Les variables {\idt largeur} et {\idt hauteur} sont déclarées dans le programme, en
dehors des procédures~; on dit qu'elles sont {\em globales}\index{Variable globale}. Une procédure
peut utiliser leur valeur sans rien dire de spécial. Par contre, pour les {\em modifier}, il faut
le déclarer au début du corps de la procédure par {\mcl global}\index{global}\ {\idt largeur}.
Sans cette déclaration, \python\ considérerait que l'affectation {\idt largeur = 0} crée une
nouvelle variable locale {\idt largeur}, sans rapport avec la variable globale, qui garderait la
valeur 0 (et {\idt largeur += 1} sans {\idt largeur = 0} avant provoquerait l'erreur
`{\tt UnboundLocalError:\ local variable 'largeur' referenced before assignment}', la variable
locale {\idt largeur} étant utilisée avant d'avoir reçu une valeur).

L'opérateur {\idt +} appliqué à deux chaînes de caractères les met bout à bout et la
fonction {\idt str} donne la chaîne de caractères qui représente un nombre (par
exemple {\idt str(25)} vaut {\cha \char34{}25\char34}). On ne peut pas mettre bout à bout une
chaîne et un nombre~: {\cha \char34{}(\char34}\ {\idt + largeur} provoque l'erreur
`{\tt TypeError:\ can only concatenate str (not "int") to str}'.

\section{Entrée d'une valeur au clavier}
On peut obtenir une valeur en la tapant au clavier de
l'ordinateur. Pour
ce faire la fonction {\idt lireEntier()} attend que l'on tape un nombre entier et rend
sa valeur. La question posée par le programme (écrite avec {\idt write}, voir ci-dessous)
s'affiche dans la barre du bas de la fenêtre du \robot, et la réponse que l'on tape au clavier
s'inscrit à sa suite (on peut aussi la taper dans le terminal d'où le programme a été lancé).
On écrit par exemple~:

\mbox{\blc  \blk\blk\idt  write({\cha \char34{}n\ =\ \char34})}\par
\mbox{\blc  \blk\blk\idt  n\ =\ lireEntier()}\par

\noindent pour donner à la variable {\idt n} la valeur tapée. On peut
corriger les fautes de frappe avec la touche  %touche d'effacement
\framebox[11.mm][l]{$\leftarrow$~~~}\hspace{0.3ex}
d'effacement. Il faut terminer l'entrée de la valeur par la frappe de la touche %retour a la ligne
\framebox[11.mm][l]{$\hookleftarrow$~~~}\hspace{0.3ex}
de retour à la ligne. Si ce que l'on a tapé n'est pas un entier, le programme le
redemande. Sur un clavier français, on peut taper les chiffres avec ou sans la touche
des majuscules.

De même {\idt lireReel()} rend un nombre rationnel (de type {\idt float}), {\idt lireChaine()}
rend la chaîne de caractères tapée (sans apostrophes ni guillemets) et {\idt lireCar()} rend
le caractère de la première touche frappée dans la fenêtre du \robot\ (sans attendre la touche de
retour à la ligne). En \pascal, une seule procédure {\idt readln} permettait de lire une valeur
de n'importe quel type. En \python, chaque fonction de lecture rend une valeur d'un type bien précis
(la fonction {\idt input} de \python\ rend toujours une chaîne de caractères, qu'il faudrait
convertir en entier avec {\idt int}~; {\idt lireEntier} le fait pour nous). Par contre certaines
fonctions, comme {\idt write}, {\idt str} ou {\idt abs}, acceptent des paramètres de types
différents~: on dit qu'elles sont {\em polymorphes}\index{Polymorphe} (alors que d'autres fonctions, comme
{\idt lireEntier} ou {\idt sqrt}, sont {\em monomorphes}\index{Monomorphe} en ce sens que leurs
paramètres et leur résultat ont toujours le même type).

L'instruction {\idt readln()} permet d'attendre que la touche %retour a la ligne
\framebox[11.mm][l]{$\hookleftarrow$~~~}\hspace{0.3ex}
de retour à la ligne soit enfoncée. Elle est donc utile pour laisser le temps de lire ce que
le programme a écrit avant la fin de son exécution.

\section{Affichage de la valeur d'une expression}
Pour afficher des textes dans le terminal (et dans la barre du bas de la fenêtre du
\robot), on
peut utiliser l'instruction {\idt write(}$t_1${\idt ,} $t_2${\idt ,} \ldots{\idt )} qui écrit 
$t_1$, $t_2$, \ldots\ les uns
à la suite des autres, sur la même ligne. L'instruction
{\idt writeln(}$t_1${\idt ,} $t_2${\idt ,} \ldots{\idt )} fait de même puis passe à la ligne suivante.
Les $t_i$ peuvent être des chaînes de caractères ou des expressions entières ou rationnelles,
dont la valeur est écrite~: si {\idt x} vaut 3, {\idt writeln(}{\cha \char34{}x = \char34}{\idt , x)}
écrit {\tt x = 3}. On peut aussi transformer soi-même la valeur d'une expression $E$ en chaîne de caractères
avec {\idt str(}$E${\idt )}, comme le fait le programme ci-dessous. Par exemple
ce programme lit
un entier $n$
puis écrit sa factorielle $n! = (1 \ast 2 \ast \ldots \ast (n - 1) \ast n)$. Pour redemander
la valeur de $n$ tant qu'elle est incorrecte, il utilise une fonction récursive {\idt lireN}
(voir le chapitre \ref{ch:rec}) qui rend la valeur lue~:
\input python/Factorielle.tex

\noindent La variable {\idt n} de la fonction {\idt lireN} est locale à cette fonction~: c'est
une autre variable que la variable {\idt n} du programme, qui reçoit la valeur rendue par
{\idt lireN()}. La valeur maximale de $n$ est fixée à 20, comme dans la version \ocaml\ de
ce programme où $21!$ dépasse le plus grand entier. En \python, les entiers n'ont pas de limite~:
on pourrait calculer exactement $100!$, qui s'écrit avec 158 chiffres.
\vfill

\noindent \rule{137mm}{0.1mm} %---------------------------------------------

\noindent\begin{exercices}{\ \\

\tiret Écrire un programme qui dessine un cadre le plus près possible du bord de la fenêtre de déplacement
du \robot.

\tiret Écrire un programme qui esquisse les dessins à  compléter de la figure \ref{fig:completer-les-dessins} 
(page \pageref{fig:completer-les-dessins}).

\tiret La numération babylonienne \cite{guitel}, qui date d'environ 2000 ans av. J.-C., est en
base 60. C'est donc une numération sexagésimale tout comme celle que nous utilisons encore
pour noter l'heure. Les unités, de 1 à 59, sont notées de manière additive en accumulant des ``clous'' (pour
1) et des ``chevrons'' (pour 10) regroupés et superposés comme indiqué sur la figure
ci-dessous.\\[3.9cm]
\noindent\includePostScript{Babyloniens}{0.95}{0.95}%[3.9cm]

\noindent Au delà de 60, la notation est positionnelle, la valeur d'un chiffre occupant
la $p^{\grave{e}me}$ position (comptée de droite à gauche à partir de zéro) étant
multipliée par $60^{p}$. L'absence d'unités sexagésimales en position médiale (que
nous no\-te\-rions par 0) fut d'abord marquée par un espace puis, vers le $3^{\grave{e}me}$
siècle av. J.-C., par ``deux clous inclinés''. Les Babyloniens ont donc mis 15 siècles pour
inventer une notation pour le
zéro sans toutefois en comprendre la valeur calculatoire$\:$! Par exemple 
le nombre $3657 = 1 \times 60^{2} + 0 \times 60^{1} + 57 \times 60^{0}$
est représenté par la juxtaposition des chiffres $1$, $0$ et 
$57 = 50 + 7$~:

\pagebreak
%\vspace{-6mm}
\mbox{}\\[1.9cm]
\noindent\includePostScript{Babyloniens3657}{0.95}{0.95}%[2.4cm]

Écrire un programme \python\ de conversion
d'un nombre po\-si\-tif quel\-con\-que dans la nu\-mé\-ra\-tion babylonienne. 

\tiret La numération égyptienne ($3^{\grave{e}me}$ 
millénaire av. J.-C.) est décimale et additive \cite{guitel}. Elle utilise des hiéroglyphes pour 
noter les
puissances de 10 (le trait vertical pour 1, l'anse pour 10, la spirale pour 100, le lotus pour 1000,
l'index pour 10~000,
le têtard pour 100 000 et le génie pour 1 000 000)~:\\[3.0cm]

\noindent\includePostScript{ChiffresEgyptiens}{0.95}{0.95}%[3.0cm]
 
\noindent La valeur d'une séquence de signes est la somme des valeurs
de ces signes~:\\[2.0cm]

\noindent\includePostScript{1234567Egyptien}{0.95}{0.95}%[3.0cm]
 
\noindent On écrit de gauche à droite ou de droite à gauche, la
lecture se faisant dans le sens où sont orientés les hiéroglyphes. 
Les signes identiques se regroupent sur une, deux ou trois lignes de
la manière suivante~:\\[0.8cm]
 
%\vspace{-\baselineskip}
\noindent\includePostScript{GroupesEgyptiens}{0.95}{0.95}%[1.8cm]
 
\noindent  Entre -2000 et -1600 ans av. J.-C., une représentation hybride (additive et partiellement
positionnelle) est apparue. Un nombre inférieur à 1000 situé avant, au dessus ou au dessous
du signe du génie du million, du têtard, de l'index ou du lotus représente le produit de ce nombre et de la
valeur de ce hiéroglyphe. Par exemple, le nombre 37 280 549 peut se décomposer en
$((372 \ast 100 000) + 80 549)$ et s'écrit~:\\[2.3cm]
 
\noindent\includePostScript{37280549Egyptien}{0.95}{0.95}
 
\noindent Écrire un programme \python\ de conversion d'un nombre dans la nu\-mé\-ra\-tion
é\-gyp\-tien\-ne.

\tiret La numération acrophonique grecque \cite{guitel} fut utilisée à partir 
du V$\grave{e}me$
siècle av. J.-C. à Athènes. C'est une numération additive décimale utilisant la base
auxiliaire 5 pour diminuer les répétitions. Les chiffres supérieurs
à un sont notés par des 
lettres (ou combinaisons de lettres) 
correspondant chacune à l'initiale (ou à l'abréviation) du nom du
nombre correspondant~:\\[1.8cm]

\noindent\includePostScript{ChiffresGrecques}{0.95}{0.95}

Écrire un programme \python\ de conversion d'un nombre in\-fé\-rieur à 100~000 dans la nu\-mé\-ra\-tion
grecque. 
Pour représenter les grands nombres, utiliser la numération acrophonique 
pour les monnaies~:

\noindent Oboles~:\\[0mm]

%\vspace{-0.5cm}
\noindent\includePostScript{ObolesGrecques}{0.95}{0.95}

\noindent Drachmes (1 drachme = 6 oboles)~:\\[0.8cm]

\noindent\includePostScript{DrachmesGrecques}{0.95}{0.95}

\noindent Talents (1 talent = 6~000 drachmes)~:\\[0.8cm]

\noindent\includePostScript{TalentsGrecques}{0.95}{0.95}

\clearpage
\noindent Par exemple le nombre 239~686~966 se représente par~:\\[1.0cm]

\noindent\includePostScript{239686966Grecque}{0.95}{0.95}

\tiret La numération romaine \cite{guitel} est additive et soustractive (tout signe placé à gauche
d'un signe de valeur supérieure s'en soustrait). Ses chiffres
sont bien connus~:\\[1.2cm]

\noindent\includePostScript{ChiffresRomains}{0.95}{0.95}

Pour les grands nombres, la barre horizontale multiplie par 1~000 la valeur du nombre qu'elle
surligne et le rectangle sans base multiplie par 100~000 la valeur du nombre qu'il enferme,
comme dans l'exemple 135~792~468 ci-dessous~:\\[1.4cm]

\noindent\includePostScript{135792468Romain}{0.95}{0.95}

Écrire un programme \python\ de conversion d'un nombre dans la nu\-mé\-ra\-tion
ro\-mai\-ne.

\tiret Écrire un programme de dessin de la courbe de Sierpinski à un ordre
quelconque de la figure \ref{fig:CourbSierp} (page \pageref{fig:CourbSierp}).

\tiret Un carré magique d'ordre $n$ est un tableau carré de $n^{2}$ nombres
tel que la somme des nombres de chaque ligne, de chaque colonne et de chaque
dia\-go\-na\-le est la même. Des exemples sont donnés à la figure \ref{fig:Cmagic}
(page \pageref{fig:Cmagic}). Notons $C_{i, j}$ le nombre figurant à la
$i^{\grave{e}me}$ ligne et la $j^{\grave{e}me}$ colonne du carré comptées à
partir de 1 et numérotées de bas en haut pour les lignes et de gauche à
droite pour les colonnes. Écrire un programme \python\ qui cons\-truit des
carrés magiques d'ordre $n$ impair en commençant par placer 1 dans
$C_{i, j} = C_{\frac{n + 1}{2}, n}$ puis affecte les nombres 2, 3, \ldots, $n^{2}$
aux éléments $C_{i, j}$ obtenus en augmentant $i$ et $j$ de 1 modulo $n$ à chaque fois
pendant $n - 1$ pas, et en diminuant $j$ de 1 tout en laissant $i$ inchangé à chaque 
$n^{\grave{e}me}$ pas.

\noindent $\Box$}\end{exercices}
\begin{figure}[htbp]
 \includevspacePostScript[9.2cm]{CourbeSierpinski0a5}{0.95}{0.95}
 \caption{Courbes de Sierpinski d'ordres 0 à 5}\label{fig:CourbSierp}
\end{figure}
\begin{figure}[htbp]
 \includevspacePostScript[6.0cm]{CarreMagique}{0.95}{0.95}
 \caption{Carrés magiques d'ordres 3, 5 et 7}\label{fig:Cmagic}
\end{figure}

\pagebreak

\begin{corrige}{
Nous donnons un programme (sans le corps de la procédure technique {\idt ecrireCentre}).
Les autres se trouvent dans le dossier\newline {\idt livre/\allowbreak 22\_variables\_E\_S}~:
\input python/CarreMagique.tex
\noindent $\Box$}\end{corrige}
\vfill

%%%%%%%%%%%%%%%%%%%%%%%%%%%%%%%%%%%%%%%%%%%%%%%%%%%%%%%%%%%%%%%%%%%%%%%%%%%
\noindent \framebox[\textwidth]{
\begin{minipage}{\myparboxwidth}
\setlength{\parindent}{2em}
\small \sf
\vspace{\abovedisplayskip}
\section*{Addenda}
Une \underline{variable}\index{Variable} sert à donner un {\em nom} à une mémoire qui conserve
une {\em valeur} représentée
selon un code déterminé par le {\em type} de la variable. Pour un codage binaire, on
peut imaginer que la mémoire est une suite de lampes éteintes (pour représenter
le bit 0) ou allumées (pour le bit 1).
Dans beaucoup de langages, comme \pascal, la valeur
initiale d'une va\-ria\-ble est {\em indéterminée} puisqu'elle dépend de l'état
de la mémoire au moment du début de l'exécution du programme, ce qui est une source
d'erreurs. En \python, une variable reçoit une valeur initiale dès sa création.
L'instruction d'\underline{affectation}\index{Affectation} permet de
conserver une nouvelle valeur en mo\-di\-fiant physiquement
l'état de cette mémoire. La nouvelle valeur est donnée par une expression (qui peut
utiliser l'ancienne valeur). La notion informatique de variable est donc très différente
de la notion mathématique homonyme. En mathématiques une variable désigne une valeur
inconnue fixée une fois pour toutes, comme les constantes d'un
programme \python\ dont on ne change jamais la valeur. Pour un mathématicien, une variable informatique
serait la suite des valeurs qu'elle prend à sa création puis après chaque
affectation.

\vspace{\baselineskip}


On peut déclarer des variables {\em locales} à une procédure (ou fonction),
en écrivant {\idt x =} $\ldots$ dans le corps de la procédure, comme
la variable {\idt n} de la fonction {\idt lireOrdre} du programme {\idt CarreMagique}.
Dans ce cas la variable n'est utilisable que dans la suite du corps de la procédure.
À l'exécution, elle est créée à chaque exécution du corps
de la procédure. Pour une procédure récursive,
il s'agit à chaque appel d'une nouvelle variable sans aucun
lien avec les instances correspondant aux appels récursifs antérieurs.

\vspace{\abovedisplayskip}
\end{minipage}
}
%%%%%%%%%%%%%%%%%%%%%%%%%%%%%%%%%%%%%%%%%%%%%%%%%%%%%%%%%%%%%%%%%%%%%%%%%%%%
%%%%%%%%%%%%%%%%%%%%%%%%%%%%%%%%%%%%%%%%%%%%%%%%%%%%%%%%%%%%%%%%%%%%%%%%%%%
\noindent \framebox[\textwidth]{
\begin{minipage}{\myparboxwidth}
\setlength{\parindent}{2em}
\small \sf
\vspace{\abovedisplayskip}
Dans le corps de la procédure on peut également utiliser des variables
{\em glo\-ba\-les} dont la déclaration figure dans le programme,
en dehors des procédures, comme {\idt largeur} et {\idt hauteur} dans le programme
{\idt CalculerTailleFenetreRobot} ou {\idt i} et {\idt j} dans le programme {\idt CarreMagique}.
Pour qu'une procédure puisse les modifier par affectation, il faut les nommer après le
mot-clé {\mcl global} au début de son corps. Ces variables globales existent
pendant toute la durée de l'exécution du programme. Par conséquent en cas
d'appels récursifs d'une procédure, il s'agit toujours de la
même variable globale.

Enfin des variables locales et globales peuvent porter le même nom (comme la variable
{\idt n} de la fonction {\idt lireOrdre} et la variable {\idt n} du programme {\idt CarreMagique}). Dans
ce cas c'est la variable locale qui est visible dans la procédure, la variable
globale continuant à exister mais étant inaccessible. En \python, une variable est locale
à une procédure dès qu'elle reçoit une valeur par affectation dans le corps de cette procédure
(sauf si elle est déclarée {\mcl global}).

\vspace{\abovedisplayskip}
\end{minipage}
}
%%%%%%%%%%%%%%%%%%%%%%%%%%%%%%%%%%%%%%%%%%%%%%%%%%%%%%%%%%%%%%%%%%%%%%%%%%%% 

